\begin{abstract}
Panels of reported cases are the standard data for evaluating non-pharmaceutical interventions, yet
difference-in-differences estimators on case counts or log cases rest on assumptions that epidemic
data violate: transmission is multiplicative and self-exciting, policies are adopted in response to
rising cases, and units infect one another. We place parallel trends on the log reproduction number
inside a renewal model whose infection pressure is imported from neighbouring units, and estimate
the resulting effects by quasi-differenced generalised method of moments. The quasi-difference
removes unrestricted unit effects while using instruments built only from the past, so adoption
that responds to the epidemic is allowed. Moment contributions form a martingale difference array
across weeks even when neighbouring clusters exchange infections, so the conventional
cluster-robust sandwich applies without spatial adjustments, and the week effects are profiled out
analytically, leaving a three-dimensional problem: on a panel of 3091 counties over 26 weeks a fit
takes 0.05 seconds and a global root search 0.73 seconds. A pre-specified simulation with 8000
replications per scenario under adoption feedback, overdispersion and small counts gives bias
between $-0.002$ and $0.001$ for the direct effect and 94.5\% to 95.4\% coverage of nominal 95\%
intervals, with correct size under the null, while two-step reproduction-number regressions and
Poisson fixed effects are biased and reject far too often. An estimand of cases averted propagates
the counterfactual through the renewal equation and the network. In an application to state mask
orders and county COVID-19 cases, the order is set above the unit, so interior counties have no
variation that separates own from neighbours' orders; the total effect of a state order on an
interior county, $-0.466$ (95\% interval $-1.091$ to $0.159$), is what the panel identifies, and
the profile of the estimating equations makes the remaining ambiguity visible.
\end{abstract}
\begin{keywords}
Causal inference; Difference-in-differences; Generalized method of moments; Infectious disease epidemiology; Interference; Renewal equation.
\end{keywords}
